% Part: lambda-calculus
% Chapter: introduction
% Section: free-variables

\documentclass[../../../include/open-logic-section]{subfiles}

\begin{document}

\olfileid{lam}{syn}{fv}
\olsection{मुक्त चरे}

लॅम्डा कलनात फलनांचा अभ्यास केला जातो, आणि लॅम्डा अमूर्तीकरणातून
फलने मिळतात. सहज समजुतीने, $\lambd[x][M]$ हे असे फलन आहे की
फलनाचा फलसाधक~$x$ ला दिल्यावर त्याची मूल्ये~$M$ ने
ठरतात. पण~$M$ मधील $x$ ची प्रत्येक आस्थिती यासाठी संबंधित नसते:
$M$ मध्ये आणखी एखादे अमूर्तीकरण $\lambd[x][N]$ असेल, तर~$N$
मधील $x$ च्या आस्थिती $\lambd[x][N]$ शी संबंधित असतात;
$\lambd[x][M]$ शी नाही. म्हणून $\lambd[x][M]$ मधील
लॅम्डा अमूर्तीकरण $\lambd[x]$ हे~$M$ मधील $x$ च्या अशा
आस्थिती \emph{बद्ध करते} ज्या दुसऱ्या लॅम्डा अमूर्तीकरणाने
आधीच बद्ध केलेल्या नाहीत---म्हणजे~$M$ मधील $x$ च्या
\emph{मुक्त} आस्थिती.

\begin{defn}[व्याप्ती]
पद~$N$ मध्ये $\lambd[x][M]$ आढळत असेल, तर त्यातील~$M$
ची संबंधित आस्थिती ही~$\lambd[x]$ ची
\emph{व्याप्ती} असते.
\end{defn}


\begin{defn}[मुक्त आणि बद्ध आस्थिती]
  पद~$M$ मधील चर $x$ ची एखादी आस्थिती $\lambd[x]$ च्या
  व्याप्तीत नसेल तर ती \emph{मुक्त} असते; अन्यथा ती
  \emph{बद्ध} असते. $\lambd[x][M]$ मधील चर~$x$ ची
  एखादी आस्थिती सुरुवातीच्या $\lambd[x]$ ने बद्ध केलेली
  असते, जर आणि फक्त जर~$M$ मधील $x$ ची ती आस्थिती मुक्त असते.
\end{defn}

\begin{ex}
  $\lambd[x][x y]$ मध्ये $x$ आणि $y$ हे दोन्ही $\lambd[x]$
  च्या व्याप्तीत आहेत; म्हणून $x$ हे $\lambd[x]$ ने बद्ध होते.
  $y$ कोणत्याही $\lambd[y]$ च्या व्याप्तीत नसल्याने ते मुक्त
  असते. $\lambd[x][x x]$ मध्ये $x$ च्या दोन्ही आस्थिती
  $\lambd[x]$ ने बद्ध होतात, कारण $xx$ मध्ये त्या दोन्ही मुक्त
  आहेत. $((\lambd[x][xx])x)$ मध्ये~$x$ ची शेवटची आस्थिती
  मुक्त आहे, कारण ती कोणत्याही $\lambd[x]$ च्या व्याप्तीत नाही.
  $\lambd[x][(\lambd[x][x])x]$ मध्ये पहिल्या $\lambd[x]$
  ची व्याप्ती $(\lambd[x][x])x$ आहे आणि दुसऱ्या $\lambd[x]$
  ची व्याप्ती शेवटून दुसरी असलेली~$x$ ची आस्थिती आहे.
  $(\lambd[x][x])x$ मध्ये $x$ ची शेवटची आस्थिती मुक्त,
  तर शेवटून दुसरी बद्ध आहे. म्हणून
  $\lambd[x][(\lambd[x][x])x]$ मधील $x$ ची शेवटून दुसरी
  आस्थिती दुसऱ्या $\lambd[x]$ ने, आणि शेवटची आस्थिती
  पहिल्या $\lambd[x]$ ने बद्ध होते.
\end{ex}

पद $P$ मधील चरांच्या सर्व आस्थिती तपासून मुक्त चरांचा संच
मिळवता येतो. हा संच $\FV{P}$ असा दर्शवतात; त्याची
स्वाभाविक व्याख्या पुढीलप्रमाणे आहे:

\begin{defn}[पदाचे मुक्त चरे] \ollabel{def:fv}
  पदाच्या \emph{मुक्त चरांचा} संच विगमनाने पुढीलप्रमाणे परिभाषित करतात:
  \begin{enumerate}
    \item $\FV{x} = \{x\}$ \ollabel{def:fv1}
    \item $\FV{\lambd[x][N]} = \FV{N} \setminus \{x\}$ \ollabel{def:fv2}
    \item $\FV{PQ} = \FV{P} \cup \FV{Q}$ \ollabel{def:fv3}
  \end{enumerate}
\end{defn}

\begin{prob}
  \begin{enumerate}
  \item या पदातील $\lambd[g]$ आणि दोन $\lambd[x]$ यांच्या
    व्याप्ती ओळखा: $\lambd[g][(\lambd[x][g (x x)]) \lambd[x][g (x x)]]$.
  \item $\lambd[g][(\lambd[x][g (x x)]) \lambd[x][g (x x)]]$
    मधील चरांच्या सर्व आस्थिती बद्ध आहेत का? असल्यास,
    त्यांना अनुक्रमे कोणती अमूर्तीकरणे बद्ध करतात?
    \item $\FV{\lambd[x][(\lambd[y][(\lambd[z][x y]) z]) y]}$
      काढा.
  \end{enumerate}
\end{prob}

\begin{explain}
मुक्त चर म्हणजे जणू बाह्य जगाचा (\emph{परिसराचा}) संदर्भ.
मुक्त चरे असलेले पद अंशतःच निश्चित झालेले पद मानता येते,
कारण त्याचे वर्तन आपण परिसर कसा ठरवतो यावर अवलंबून असते.
उदाहरणार्थ, $\lambd[x][f x]$ हे पद फलसाधक~$x$ स्वीकारून
त्यावर $f$ लावून मिळणारा परिणाम देते; येथे चर~$f$ मुक्त आहे.
या पदाचे मूल्य ते ज्या परिसरात आहे त्यावर, विशेषतः त्या
परिसरातील $f$ च्या मूल्यावर, अवलंबून असते.

या पदावर अमूर्तीकरण केल्यास $\lambd[f][\lambd[x][f
    x]]$ मिळते. आता हे पद परिसरातील चर $f$ वर अवलंबून
नसते, कारण ते दोन फलसाधक स्वीकारून पहिला दुसऱ्यावर
लावल्याचा परिणाम देणारे फलन दर्शवते. परिसरातील $f$
बदलला तरी या पदाच्या वर्तनावर परिणाम होत नाही: हे पद
फक्त फलसाधक म्हणून दिलेले मूल्य वापरते, परिसरातील
$f$ चे मूल्य वापरत नाही.
\end{explain}

\begin{defn}[बंद पद, संयोजक]
  मुक्त चरे नसलेल्या पदाला \emph{बंद पद} किंवा
  \emph{संयोजक} म्हणतात.
\end{defn}

\begin{lem}\ollabel{lem:fv}
  \begin{enumerate}
    \item \ollabel{lem:fv-abs} $y \neq x$ असेल, तर $y \in
      \FV{\lambd[x][N]}$ जर आणि फक्त जर $y \in \FV{N}$.
    \item \ollabel{lem:fv-app} $y \in \FV{PQ}$ जर आणि
      फक्त जर $y \in \FV{P}$ किंवा $y
      \in \FV{Q}$.
    \end{enumerate}
\end{lem}

\begin{proof}
  सरावासाठी.
\end{proof}

\begin{prob}
  \olref[lam][syn][fv]{lem:fv} सिद्ध करा.
\end{prob}

\end{document}
